Notation as in the statement and in Theorem \claimref{synthesis.incidence-vacua-thermofield-real-capacity}, with $X(x) = \nO + x - \tfrac12 Q(x)\ninf$ and $L_{x}$ as in Proposition \claimref{one-relation.spacetime-manifold-vacuum-lines}; the computations are in the eight-dimensional Fock module of $\Cl(2,4)$, where the generators of positive square are Hermitian and those of negative square anti-Hermitian.

(i) $(e_{1}e_{2})^{2} = -e_{1}^{2}e_{2}^{2} = 1$, so $S_{\lambda} = \cosh(\lambda/2) + \sinh(\lambda/2)e_{1}e_{2}$ (\claimref{std-algebra.bivector-exponential}) with inverse $\cosh(\lambda/2) - \sinh(\lambda/2)e_{1}e_{2}$. Since $(e_{1}e_{2})^{*} = e_{2}^{*}e_{1}^{*} = (-e_{2})e_{1} = e_{1}e_{2}$, $S_{\lambda}^{*} = S_{\lambda} \ne S_{\lambda}^{-1}$ for $\lambda \ne 0$: the boost is not unitary. The bivector $e_{1}e_{2}$ commutes with $e_{\pm}$, hence with $\nO$, $\ninf$ and $N$. On the plane, $S_{\lambda}(e_{1} \pm e_{2})S_{\lambda}^{-1} = e^{\mp\lambda}(e_{1} \pm e_{2})$, so $x = \tfrac12 v(e_{1} + e_{2}) + \tfrac12 u(e_{1} - e_{2})$ goes to the point with null coordinates $e^{\lambda}u$, $e^{-\lambda}v$, and $Q$ is unchanged. As $S_{\lambda}$ fixes $\nO$ and $\ninf$, $S_{\lambda}X(x)S_{\lambda}^{-1} = \nO + x_{\lambda} - \tfrac12 Q(x)\ninf = X(x_{\lambda})$, and $S_{\lambda}$, which commutes with the chirality, maps the kernel of $X(x)$ on a half onto that of $X(x_{\lambda})$. The script checks these identities for three rational rotors ($\cosh(\lambda/2), \sinh(\lambda/2)$ on the hyperbola) and four points.

(ii) $\Omega_{x}$ and $\Psi_{x}$ lie in different spectator sectors ($\ket{0}_{+}$ and $\ket{0}_{-}$ differ in the occupation of $c$), which $N$ preserves, so the cross terms vanish. $\Omega_{x} = \ket{0}_{+} - u\,a^{*}b^{*}\ket{0}_{+}$ is a sum of orthogonal number states with $N = 0, 1$, so $\|\Omega_{x}\|^{2} = 1 + u^{2}$ and $\langle\Omega_{x}, N\Omega_{x}\rangle = u^{2}$; likewise $\|\Psi_{x}\|^{2} = 1 + v^{2}$, $\langle\Psi_{x}, N\Psi_{x}\rangle = v^{2}$. The formula follows; the script checks it at six weight pairs including complex ones.

(iii) Along the orbit, $u_{\lambda}^{2}v_{\lambda}^{2} = Q^{2}$, so $\min(u_{\lambda}^{2}, v_{\lambda}^{2}) \le |Q| \le \max(u_{\lambda}^{2}, v_{\lambda}^{2})$ and, $f$ being increasing, $f(|Q|)$ lies in every segment (every value of which is attained, Theorem \claimref{synthesis.incidence-vacua-thermofield-real-capacity}(iii), equivalently (ii) here). If $Q \ne 0$, the orbit point with $e^{2\lambda} = |v/u|$ has $u_{\lambda}^{2} = v_{\lambda}^{2} = |Q|$, where the segment is the single value $f(|Q|)$. If $Q = 0$ and $x \ne 0$, say $u = 0 \ne v$, the segments are $[0, f(e^{-2\lambda}v^{2})]$, whose intersection over $\lambda$ is $\{0\}$; at $x = 0$ every vacuum is the ascent vacuum. So an orbit-constant ascent occupation equals $f(|Q(x)|)$. The ascent marginal of a vacuum on $L_{x}$ is diagonal (Theorem \claimref{synthesis.incidence-vacua-thermofield-real-capacity}(iii), with Theorem \claimref{singularity.descended-form-state-independent}), so it is the Gibbs state with $e^{-\beta} = \langle N\rangle/(1 - \langle N\rangle) = |Q|$. Existence: the choice in (iv) (any vacuum where $|u| = |v|$) gives a section with this marginal at every point. The script checks, on three orbits and six boost factors each, that every segment contains $f(|Q|)$ strictly inside except at the collapse point, where both endpoints equal it, with pinned values $9/13$, $4/229$, $1/5$.

(iv) With $A = |\lambda_{1}|^{2}$, $B = |\mu|^{2}$, the equation $\langle N\rangle = f(|uv|)$ reads, after clearing denominators, $A(u^{2} - |uv|) + B(v^{2} - |uv|) = 0$, that is $A|u|(|u| - |v|) = B|v|(|u| - |v|)$; for $|u| \ne |v|$ this is $A|u| = B|v|$, which forces $AB \ne 0$ when $uv \ne 0$ (on the light cone the solution is the ascent vacuum, $\Omega_{x}$ at $u = 0$). By Theorem \claimref{synthesis.incidence-vacua-thermofield-real-capacity}(iii) such a vacuum has a mixed $(a,b)$ state; its purity is that of the mixture of the two pure $(a,b)$ states of $\Omega_{x}$ and $\Psi_{x}$ (the spectator distinguishes the halves) with weights proportional to $A(1 + u^{2})$ and $B(1 + v^{2})$, that is to $w_{1} = |v|(1 + u^{2})$ and $w_{2} = |u|(1 + v^{2})$, the two $(a,b)$ states being orthogonal in the number basis of the pair. The script checks the occupation, the diagonal marginal and the purity at three points with rational square roots, with pinned purities $130993/231361$, $100521361/195132961$, $70673/105625$, and that neither spanning vacuum has the boost-constant occupation there.

(v) From (iii), $\beta(x) = -\ln|Q(x)|$, which is constant on a boost orbit and strictly monotone in $|Q|$ across orbits; it equals $2\pi$ iff $|Q(x)| = e^{-2\pi}$.

\emph{Remarks (argued, not part of the claim).} The value $2\pi$ that Premise \claimref{singularity.r-regularity-fixed-locus} assigns in rapidity normalisation is the same on every Killing orbit; the boost-constant incidence reading instead gives $-\ln|Q|$, so the two agree on one hyperboloid only, and only after the scale fixed by the Hermitian structure. Because the boost is not unitary, stationarity here means constancy of the ascent marginal along the orbit, not invariance of a state under a unitary flow. A dilation $x \mapsto e^{\mu}x$ would shift $\beta$ by $-2\mu$; this inverse temperature behaves as a dilation coordinate, which is where the incidence reading and (R) part company.
